\documentclass{beamer}
% =====================================================================
% finance-beamer.tex — shared Beamer style for the LLM-in-Finance course
% =====================================================================
% Reusable slide style. Each deck \input's this immediately after
% \documentclass{beamer}. It provides:
%   * the Madrid theme + book colour palette (matches the printed book)
%   * two book-style framing blocks:
%       \begin{bigpicture} ... \end{bigpicture}   ("the bigger picture")
%       \begin{underhood}   ... \end{underhood}    ("under the hood")
%     These mirror the book's `context` and `deepdive` tcolorboxes so the
%     same intuition-vs-mechanism scaffold the reader meets in the book
%     reappears on the slides.
%   * a Python listings style for code frames
%   * appendix conventions: \appendixstart drops a divider and switches
%     to non-numbered backup frames so "extra / less central" material
%     lives after the main talk without inflating the slide count.
%
% Usage:
%   \documentclass{beamer}
%   % =====================================================================
% finance-beamer.tex — shared Beamer style for the LLM-in-Finance course
% =====================================================================
% Reusable slide style. Each deck \input's this immediately after
% \documentclass{beamer}. It provides:
%   * the Madrid theme + book colour palette (matches the printed book)
%   * two book-style framing blocks:
%       \begin{bigpicture} ... \end{bigpicture}   ("the bigger picture")
%       \begin{underhood}   ... \end{underhood}    ("under the hood")
%     These mirror the book's `context` and `deepdive` tcolorboxes so the
%     same intuition-vs-mechanism scaffold the reader meets in the book
%     reappears on the slides.
%   * a Python listings style for code frames
%   * appendix conventions: \appendixstart drops a divider and switches
%     to non-numbered backup frames so "extra / less central" material
%     lives after the main talk without inflating the slide count.
%
% Usage:
%   \documentclass{beamer}
%   % =====================================================================
% finance-beamer.tex — shared Beamer style for the LLM-in-Finance course
% =====================================================================
% Reusable slide style. Each deck \input's this immediately after
% \documentclass{beamer}. It provides:
%   * the Madrid theme + book colour palette (matches the printed book)
%   * two book-style framing blocks:
%       \begin{bigpicture} ... \end{bigpicture}   ("the bigger picture")
%       \begin{underhood}   ... \end{underhood}    ("under the hood")
%     These mirror the book's `context` and `deepdive` tcolorboxes so the
%     same intuition-vs-mechanism scaffold the reader meets in the book
%     reappears on the slides.
%   * a Python listings style for code frames
%   * appendix conventions: \appendixstart drops a divider and switches
%     to non-numbered backup frames so "extra / less central" material
%     lives after the main talk without inflating the slide count.
%
% Usage:
%   \documentclass{beamer}
%   \input{../_style/finance-beamer.tex}
%   \title{...}\subtitle{...}\author{...}\institute{...}\date{\today}
%   \begin{document} ... \end{document}
% =====================================================================

\usetheme{Madrid}
\usecolortheme{whale}
\setbeamertemplate{navigation symbols}{}
\setbeamertemplate{caption}[numbered]
\setbeamercovered{transparent}

% --- Density controls: keep dense, equation-heavy frames on one slide ---
% Slightly smaller body text + tighter lists/blocks. Frame titles and the
% Madrid chrome keep their normal size, so the deck still reads as Madrid,
% just with more vertical room for content.
\setbeamerfont{normal text}{size=\small}
\setbeamertemplate{itemize/enumerate body begin}{\small}
\setbeamertemplate{itemize/enumerate subbody begin}{\footnotesize}
% Tighter block spacing.
\setbeamertemplate{block begin}{%
  \par\vskip2pt%
  \begin{beamercolorbox}[colsep*=.75ex,leftskip=1ex,rightskip=1ex]{block title}%
    \usebeamerfont*{block title}\insertblocktitle%
  \end{beamercolorbox}%
  {\parskip0pt\vskip-1pt}%
  \begin{beamercolorbox}[colsep*=.75ex,vmode,leftskip=1ex,rightskip=1ex]{block body}%
  \vskip1pt}
\setbeamertemplate{block end}{\end{beamercolorbox}\vskip2pt}

\usepackage{amsmath, amssymb}
\usepackage{booktabs}
\usepackage{xcolor}
\usepackage{listings}
\usepackage[skins, breakable]{tcolorbox}

% --- Book colour palette (kept in sync with book/preamble.tex) ---
\definecolor{linkblue}{RGB}{14, 75, 140}
\definecolor{deepdivebg}{RGB}{235, 244, 255}
\definecolor{narrativebg}{RGB}{255, 252, 240}
\definecolor{narrativerule}{RGB}{180, 130, 40}
\definecolor{steelblue}{RGB}{70, 130, 180}
\definecolor{forestgreen}{RGB}{34, 139, 34}
\definecolor{crimson}{RGB}{220, 20, 60}
\definecolor{darkorange}{RGB}{255, 140, 0}
\definecolor{codebg}{RGB}{248, 248, 248}
\definecolor{coderule}{RGB}{210, 210, 210}
\definecolor{keywordblue}{RGB}{0, 100, 180}
\definecolor{commentgray}{RGB}{120, 120, 120}
\definecolor{stringgreen}{RGB}{30, 130, 30}

% Tie the Madrid structural colour to the book's link blue.
\setbeamercolor{structure}{fg=linkblue}
\setbeamercolor{title}{fg=white}
\setbeamercolor{frametitle}{fg=white}

% --- "the bigger picture" : narrative / intuition framing -----------
\newtcolorbox{bigpicture}{%
  enhanced, breakable, boxsep=2pt,
  colback  = narrativebg,
  colframe = narrativerule,
  boxrule  = 0pt, toprule = 1.4pt, bottomrule = 0.4pt,
  leftrule = 0pt, rightrule = 0pt, arc = 0pt,
  left = 6pt, right = 6pt, top = 4pt, bottom = 5pt,
  fonttitle = \footnotesize\scshape\color{narrativerule},
  title = {the bigger picture}, attach title to upper = {\par\smallskip},
}

% --- "under the hood" : the mechanism / technical detail ------------
\newtcolorbox{underhood}{%
  enhanced, breakable, boxsep=2pt,
  colback  = deepdivebg,
  colframe = linkblue,
  boxrule  = 0pt, toprule = 1.4pt, bottomrule = 0.4pt,
  leftrule = 0pt, rightrule = 0pt, arc = 0pt,
  left = 6pt, right = 6pt, top = 4pt, bottom = 5pt,
  fonttitle = \footnotesize\scshape\color{linkblue},
  title = {under the hood}, attach title to upper = {\par\smallskip},
}

% --- Python code style ---------------------------------------------
\lstset{
  basicstyle    = \ttfamily\footnotesize,
  keywordstyle  = \color{keywordblue}\bfseries,
  commentstyle  = \color{commentgray}\itshape,
  stringstyle   = \color{stringgreen},
  showstringspaces = false,
  language      = Python,
  breaklines    = true,
  backgroundcolor = \color{codebg},
  frame         = single,
  rulecolor     = \color{coderule},
  numbers       = none,
  aboveskip     = 4pt, belowskip = 4pt,
}

% --- Appendix conventions ------------------------------------------
% \appendixstart{Title} : begins the "extra material" appendix. Frames
% after it are not counted toward the main deck's frame total, keeping
% the core talk lean while preserving the deeper / optional content.
\newcommand{\appendixstart}[1]{%
  \appendix
  \setbeamertemplate{frametitle continuation}{}%
  \begin{frame}[noframenumbering]
    \begin{center}
      {\Large\scshape\color{linkblue} Appendix}\\[4pt]
      {\large #1}\\[10pt]
      {\footnotesize Extra / optional material — beyond the core lecture.}
    \end{center}
  \end{frame}
}
% Use \begin{frame}[noframenumbering]{...} for each appendix frame.

%   \title{...}\subtitle{...}\author{...}\institute{...}\date{\today}
%   \begin{document} ... \end{document}
% =====================================================================

\usetheme{Madrid}
\usecolortheme{whale}
\setbeamertemplate{navigation symbols}{}
\setbeamertemplate{caption}[numbered]
\setbeamercovered{transparent}

% --- Density controls: keep dense, equation-heavy frames on one slide ---
% Slightly smaller body text + tighter lists/blocks. Frame titles and the
% Madrid chrome keep their normal size, so the deck still reads as Madrid,
% just with more vertical room for content.
\setbeamerfont{normal text}{size=\small}
\setbeamertemplate{itemize/enumerate body begin}{\small}
\setbeamertemplate{itemize/enumerate subbody begin}{\footnotesize}
% Tighter block spacing.
\setbeamertemplate{block begin}{%
  \par\vskip2pt%
  \begin{beamercolorbox}[colsep*=.75ex,leftskip=1ex,rightskip=1ex]{block title}%
    \usebeamerfont*{block title}\insertblocktitle%
  \end{beamercolorbox}%
  {\parskip0pt\vskip-1pt}%
  \begin{beamercolorbox}[colsep*=.75ex,vmode,leftskip=1ex,rightskip=1ex]{block body}%
  \vskip1pt}
\setbeamertemplate{block end}{\end{beamercolorbox}\vskip2pt}

\usepackage{amsmath, amssymb}
\usepackage{booktabs}
\usepackage{xcolor}
\usepackage{listings}
\usepackage[skins, breakable]{tcolorbox}

% --- Book colour palette (kept in sync with book/preamble.tex) ---
\definecolor{linkblue}{RGB}{14, 75, 140}
\definecolor{deepdivebg}{RGB}{235, 244, 255}
\definecolor{narrativebg}{RGB}{255, 252, 240}
\definecolor{narrativerule}{RGB}{180, 130, 40}
\definecolor{steelblue}{RGB}{70, 130, 180}
\definecolor{forestgreen}{RGB}{34, 139, 34}
\definecolor{crimson}{RGB}{220, 20, 60}
\definecolor{darkorange}{RGB}{255, 140, 0}
\definecolor{codebg}{RGB}{248, 248, 248}
\definecolor{coderule}{RGB}{210, 210, 210}
\definecolor{keywordblue}{RGB}{0, 100, 180}
\definecolor{commentgray}{RGB}{120, 120, 120}
\definecolor{stringgreen}{RGB}{30, 130, 30}

% Tie the Madrid structural colour to the book's link blue.
\setbeamercolor{structure}{fg=linkblue}
\setbeamercolor{title}{fg=white}
\setbeamercolor{frametitle}{fg=white}

% --- "the bigger picture" : narrative / intuition framing -----------
\newtcolorbox{bigpicture}{%
  enhanced, breakable, boxsep=2pt,
  colback  = narrativebg,
  colframe = narrativerule,
  boxrule  = 0pt, toprule = 1.4pt, bottomrule = 0.4pt,
  leftrule = 0pt, rightrule = 0pt, arc = 0pt,
  left = 6pt, right = 6pt, top = 4pt, bottom = 5pt,
  fonttitle = \footnotesize\scshape\color{narrativerule},
  title = {the bigger picture}, attach title to upper = {\par\smallskip},
}

% --- "under the hood" : the mechanism / technical detail ------------
\newtcolorbox{underhood}{%
  enhanced, breakable, boxsep=2pt,
  colback  = deepdivebg,
  colframe = linkblue,
  boxrule  = 0pt, toprule = 1.4pt, bottomrule = 0.4pt,
  leftrule = 0pt, rightrule = 0pt, arc = 0pt,
  left = 6pt, right = 6pt, top = 4pt, bottom = 5pt,
  fonttitle = \footnotesize\scshape\color{linkblue},
  title = {under the hood}, attach title to upper = {\par\smallskip},
}

% --- Python code style ---------------------------------------------
\lstset{
  basicstyle    = \ttfamily\footnotesize,
  keywordstyle  = \color{keywordblue}\bfseries,
  commentstyle  = \color{commentgray}\itshape,
  stringstyle   = \color{stringgreen},
  showstringspaces = false,
  language      = Python,
  breaklines    = true,
  backgroundcolor = \color{codebg},
  frame         = single,
  rulecolor     = \color{coderule},
  numbers       = none,
  aboveskip     = 4pt, belowskip = 4pt,
}

% --- Appendix conventions ------------------------------------------
% \appendixstart{Title} : begins the "extra material" appendix. Frames
% after it are not counted toward the main deck's frame total, keeping
% the core talk lean while preserving the deeper / optional content.
\newcommand{\appendixstart}[1]{%
  \appendix
  \setbeamertemplate{frametitle continuation}{}%
  \begin{frame}[noframenumbering]
    \begin{center}
      {\Large\scshape\color{linkblue} Appendix}\\[4pt]
      {\large #1}\\[10pt]
      {\footnotesize Extra / optional material — beyond the core lecture.}
    \end{center}
  \end{frame}
}
% Use \begin{frame}[noframenumbering]{...} for each appendix frame.

%   \title{...}\subtitle{...}\author{...}\institute{...}\date{\today}
%   \begin{document} ... \end{document}
% =====================================================================

\usetheme{Madrid}
\usecolortheme{whale}
\setbeamertemplate{navigation symbols}{}
\setbeamertemplate{caption}[numbered]
\setbeamercovered{transparent}

% --- Density controls: keep dense, equation-heavy frames on one slide ---
% Slightly smaller body text + tighter lists/blocks. Frame titles and the
% Madrid chrome keep their normal size, so the deck still reads as Madrid,
% just with more vertical room for content.
\setbeamerfont{normal text}{size=\small}
\setbeamertemplate{itemize/enumerate body begin}{\small}
\setbeamertemplate{itemize/enumerate subbody begin}{\footnotesize}
% Tighter block spacing.
\setbeamertemplate{block begin}{%
  \par\vskip2pt%
  \begin{beamercolorbox}[colsep*=.75ex,leftskip=1ex,rightskip=1ex]{block title}%
    \usebeamerfont*{block title}\insertblocktitle%
  \end{beamercolorbox}%
  {\parskip0pt\vskip-1pt}%
  \begin{beamercolorbox}[colsep*=.75ex,vmode,leftskip=1ex,rightskip=1ex]{block body}%
  \vskip1pt}
\setbeamertemplate{block end}{\end{beamercolorbox}\vskip2pt}

\usepackage{amsmath, amssymb}
\usepackage{booktabs}
\usepackage{xcolor}
\usepackage{listings}
\usepackage[skins, breakable]{tcolorbox}

% --- Book colour palette (kept in sync with book/preamble.tex) ---
\definecolor{linkblue}{RGB}{14, 75, 140}
\definecolor{deepdivebg}{RGB}{235, 244, 255}
\definecolor{narrativebg}{RGB}{255, 252, 240}
\definecolor{narrativerule}{RGB}{180, 130, 40}
\definecolor{steelblue}{RGB}{70, 130, 180}
\definecolor{forestgreen}{RGB}{34, 139, 34}
\definecolor{crimson}{RGB}{220, 20, 60}
\definecolor{darkorange}{RGB}{255, 140, 0}
\definecolor{codebg}{RGB}{248, 248, 248}
\definecolor{coderule}{RGB}{210, 210, 210}
\definecolor{keywordblue}{RGB}{0, 100, 180}
\definecolor{commentgray}{RGB}{120, 120, 120}
\definecolor{stringgreen}{RGB}{30, 130, 30}

% Tie the Madrid structural colour to the book's link blue.
\setbeamercolor{structure}{fg=linkblue}
\setbeamercolor{title}{fg=white}
\setbeamercolor{frametitle}{fg=white}

% --- "the bigger picture" : narrative / intuition framing -----------
\newtcolorbox{bigpicture}{%
  enhanced, breakable, boxsep=2pt,
  colback  = narrativebg,
  colframe = narrativerule,
  boxrule  = 0pt, toprule = 1.4pt, bottomrule = 0.4pt,
  leftrule = 0pt, rightrule = 0pt, arc = 0pt,
  left = 6pt, right = 6pt, top = 4pt, bottom = 5pt,
  fonttitle = \footnotesize\scshape\color{narrativerule},
  title = {the bigger picture}, attach title to upper = {\par\smallskip},
}

% --- "under the hood" : the mechanism / technical detail ------------
\newtcolorbox{underhood}{%
  enhanced, breakable, boxsep=2pt,
  colback  = deepdivebg,
  colframe = linkblue,
  boxrule  = 0pt, toprule = 1.4pt, bottomrule = 0.4pt,
  leftrule = 0pt, rightrule = 0pt, arc = 0pt,
  left = 6pt, right = 6pt, top = 4pt, bottom = 5pt,
  fonttitle = \footnotesize\scshape\color{linkblue},
  title = {under the hood}, attach title to upper = {\par\smallskip},
}

% --- Python code style ---------------------------------------------
\lstset{
  basicstyle    = \ttfamily\footnotesize,
  keywordstyle  = \color{keywordblue}\bfseries,
  commentstyle  = \color{commentgray}\itshape,
  stringstyle   = \color{stringgreen},
  showstringspaces = false,
  language      = Python,
  breaklines    = true,
  backgroundcolor = \color{codebg},
  frame         = single,
  rulecolor     = \color{coderule},
  numbers       = none,
  aboveskip     = 4pt, belowskip = 4pt,
}

% --- Appendix conventions ------------------------------------------
% \appendixstart{Title} : begins the "extra material" appendix. Frames
% after it are not counted toward the main deck's frame total, keeping
% the core talk lean while preserving the deeper / optional content.
\newcommand{\appendixstart}[1]{%
  \appendix
  \setbeamertemplate{frametitle continuation}{}%
  \begin{frame}[noframenumbering]
    \begin{center}
      {\Large\scshape\color{linkblue} Appendix}\\[4pt]
      {\large #1}\\[10pt]
      {\footnotesize Extra / optional material — beyond the core lecture.}
    \end{center}
  \end{frame}
}
% Use \begin{frame}[noframenumbering]{...} for each appendix frame.


\title{LLMs and Privacy}
\subtitle{Lecture 15 --- Local Deployments and Text De-identification}
\author{Juan F.\ Imbet}
\institute{Paris Dauphine -- PSL University}
\date{\today}

\begin{document}

\begin{frame}\titlepage\end{frame}

\begin{frame}{Outline}\tableofcontents\end{frame}

% =====================================================================
\section{The Privacy Imperative in Financial AI}
% =====================================================================

\begin{frame}{Why Privacy Is the Binding Constraint}
\begin{bigpicture}
January 2023: a major bank bans staff from pasting client data into ChatGPT after
a leaked memo reveals loan applications, earnings transcripts, and merger term
sheets crossing a public API. Weeks later, Samsung source code submitted to
ChatGPT is retained for training. \textbf{Once data crosses the API boundary,
there is no technical mechanism for recall.}
\end{bigpicture}

\smallskip
LLMs are useful \emph{because} they absorbed vast text --- but absorption plus cloud
inference creates a privacy surface finance is poorly equipped to govern.

\smallskip
\textbf{Four complementary engineering strategies} (the spine of this lecture):
\begin{enumerate}
  \item Local / self-hosted deployment --- remove the external boundary
  \item Text anonymisation --- strip sensitive content before the model sees it
  \item Privacy-preserving training (DP, federated learning) --- bound memorisation
  \item Privacy-aware architecture (RAG, audit, access control) --- govern the system
\end{enumerate}
\end{frame}

\begin{frame}{What Data Do LLMs Expose? Two Channels}
\begin{columns}[T]
\begin{column}{0.5\textwidth}
\begin{block}{Training channel}
Information in the pre-training / fine-tuning corpus is reconstructable from
\emph{parameters}.
\begin{itemize}
  \item Persists after the session ends
  \item Any user may recover fragments
  \item Fine-tune on proprietary docs $\Rightarrow$ permanent artefact
\end{itemize}
\end{block}
\end{column}
\begin{column}{0.5\textwidth}
\begin{block}{Inference channel}
Data in \emph{prompts} may be stored, logged, or trained on by the provider.
\begin{itemize}
  \item Operational; depends on the contract
  \item Eliminated by on-premise deployment
  \item Encountered first --- no skill needed to exploit it
\end{itemize}
\end{block}
\end{column}
\end{columns}

\medskip
\begin{alertblock}{Memorisation is real and quantifiable}
Carlini et al.\ (2021): a 1.5B-param GPT-2 memorises verbatim sequences --- emails,
phone numbers, identifiers seen \emph{once}. Larger models memorise more; repeated
sequences memorised most. Carlini et al.\ (2022) formalise a per-example
\emph{memorisation score}.
\end{alertblock}
\end{frame}

\begin{frame}{Definition: Memorisation}
\begin{block}{Memorisation (Def.\ 15.1)}
A language model $\mathcal{M}$ \emph{memorises} a string $s$ if there exists a
prompt $p$ such that $\mathcal{M}(p)$ generates $s$ with probability above a
threshold $\tau \in (0,1]$.
\end{block}

\medskip
\begin{itemize}
  \item \textbf{Exact} memorisation: $s$ reproduced verbatim.
  \item \textbf{Approximate} memorisation: bounded Levenshtein edit distance
        (min.\ single-character insert/delete/substitute edits).
\end{itemize}

\medskip
\begin{underhood}
Why rare sequences get memorised: the next-token objective rewards
\emph{generalisable} features for common patterns, but for a rare pattern the only
path to low loss is to store the specific sequence. A client name appearing once in
a fine-tuning corpus is \emph{memorised}; the word ``the'' is \emph{learned}. The
distinction is, in part, a function of frequency.
\end{underhood}
\end{frame}

\begin{frame}{Regulatory Landscape: Four Frameworks at Once}
\small
\begin{tabular}{p{1.8cm}p{3.0cm}p{5.3cm}}
\toprule
\textbf{Framework} & \textbf{Scope} & \textbf{Key bite for LLMs} \\
\midrule
GDPR (2016/679) & Personal data of EEA individuals (broadest) &
  Art.\ 22 (solely automated decisions), Art.\ 17 (erasure),
  Art.\ 83 fines: higher of EUR\,20m / 4\% global turnover \\
\addlinespace
CCPA / CPRA & California residents &
  Right to opt out of \emph{sale}; may cover sharing with model providers \\
\addlinespace
MiFID II (2014/65) & Investment services &
  Record-keeping: 5 yr (services), 7 yr (orders), retrievable \\
\addlinespace
DORA (2022/2554) & ICT third parties (since Jan 2025) &
  Critical-provider oversight; data location, audit rights, exit strategy \\
\bottomrule
\end{tabular}

\medskip
\normalsize
\begin{alertblock}{Compliance $\neq$ privacy}
A GDPR-compliant cloud provider can still create client privacy risk; a technically
private local model can still fail MiFID II record-keeping. \textbf{Assess the
frameworks together, against each use case.}
\end{alertblock}
\end{frame}

% =====================================================================
\section{Privacy Risks in LLM Deployment}
% =====================================================================

\begin{frame}{The Three Risk Mechanisms}
\begin{bigpicture}
Privacy risks map onto the \emph{confidentiality} axis of the CIA triad --- but
through mechanisms unlike a database breach. \textbf{No firewall stops membership
inference; no encryption stops an LLM from reciting a training example.}
Understanding the mechanism is a prerequisite for choosing the mitigation.
\end{bigpicture}

\medskip
\begin{enumerate}
  \item \textbf{Training-data memorisation \& leakage} --- extraction attacks recover
        rare / repeated sequences from parameters.
  \item \textbf{Inference attacks} --- membership and property inference leak
        information \emph{without} verbatim memorisation.
  \item \textbf{Prompt injection \& exfiltration} --- adversarial instructions in
        retrieved / ingested content hijack the model.
\end{enumerate}
\end{frame}

\begin{frame}{Memorisation and Extraction Attacks}
Neural LMs compress training data into parameters --- but compression is lossy only
\emph{in aggregate}. Unusual, repeated, or long examples are retained with high fidelity.

\medskip
\begin{underhood}
\textbf{Likelihood-ratio extraction} (Carlini et al., 2021):
\begin{enumerate}
  \item Generate many long continuations from varied prompts.
  \item Rank by $\dfrac{\Pr_{\text{target}}(s)}{\Pr_{\text{reference}}(s)}$ --- target
        vs.\ a smaller reference model.
  \item High ratio $\Rightarrow$ disproportionately likely to be training data
        (the reference has not seen the specific sequence).
\end{enumerate}
Recovered: verbatim emails, phone numbers, names, full PII-bearing examples.
\end{underhood}
\end{frame}

\begin{frame}{Memorisation: Why It Is Amplified in Finance}
\begin{alertblock}{Amplified in finance}
Fine-tuning on internal research / deal docs / client comms is attractive --- and
embeds those docs in parameters. A competitor querying the model can recover
fragments \emph{even behind an access-controlled API}.
\end{alertblock}

\medskip
The extraction risk scales with how \emph{rare} and how \emph{repeated} a sequence
is in the fine-tuning corpus. Proprietary deal documents are both: they appear
verbatim across drafts, and their exact phrasing exists nowhere in the public
pre-training data --- maximising the likelihood ratio that exposes them.
\end{frame}

\begin{frame}{Inference Attacks: Membership and Property}
\begin{block}{Membership Inference Advantage (Def.\ 15.2)}
For model $\mathcal{M}$ trained on $D$ and adversary $\mathcal{A}$ observing a
candidate $s$:
\[
\mathrm{Adv}(\mathcal{A}) = \bigl|\Pr[\mathcal{A}(\mathcal{M}, s){=}1 \mid s \in D]
  - \Pr[\mathcal{A}(\mathcal{M}, s){=}1 \mid s \notin D]\bigr|
\]
$0$ = no discrimination; $1$ = perfect discrimination. Basis: models assign higher
likelihood to training examples (Shokri et al., 2017).
\end{block}

\medskip
\textbf{Financial consequence.} A bank fine-tunes on merger targets; an adversary
queries each candidate name and reconstructs the \textbf{deal pipeline} with
above-chance accuracy. Same logic leaks whether a specific credit file was in training.

\medskip
\textbf{Property inference} extends this to \emph{aggregates} (mean loan size,
fraction delinquent) --- harder to detect, requires no example-level memorisation.
\end{frame}

\begin{frame}{Prompt Injection and Data Exfiltration}
\textbf{Prompt injection} (Perez \& Ribeiro, 2022): malicious instructions embedded
in processed content override system-level instructions. In RAG, retrieved documents
are concatenated into the prompt --- whoever controls a document controls instructions.

\medskip
\begin{block}{Two finance-specific exposures}
\begin{itemize}
  \item \textbf{Adverse-media screening}: an attacker publishes an article with an
        injection payload $\Rightarrow$ forces a favourable risk score or leaks the
        screening criteria.
  \item \textbf{Counterparty document ingestion}: loan applications, due-diligence
        questionnaires, ISDA agreements carrying embedded instructions.
\end{itemize}
\end{block}

\medskip
\textbf{Mitigations} (no complete solution yet):
structural separation of trusted vs.\ untrusted content $\cdot$ output sanitisation
$\cdot$ rule-based / constitutional classifiers scanning output before delivery.
\end{frame}

% =====================================================================
\section{Local and Self-Hosted LLM Deployments}
% =====================================================================

\begin{frame}{Keep Inference Inside the Perimeter}
\begin{bigpicture}
The simplest way to kill inference-time exfiltration: keep inference on the
institution's own infrastructure. For many financial workflows the mix of
regulation, contracts, and competitive sensitivity makes on-premise not merely
preferable but \emph{mandatory}. Capable open-weight models (2023--2024) made this
practically feasible for the first time.
\end{bigpicture}

\medskip
\textbf{Open-weight} = parameters publicly released; inference needs no third-party
API call. (Distinct from fully open-source --- weights yes, full training
code / data often no.) What matters for privacy: \textbf{no data crosses an external
boundary.}
\end{frame}

\begin{frame}{Open-Weight Models: Llama, Mistral, Phi}
\small
\begin{tabular}{p{2.1cm}p{1.5cm}p{2.4cm}p{4.0cm}}
\toprule
\textbf{Family} & \textbf{Params} & \textbf{4-bit footprint} & \textbf{Niche} \\
\midrule
Llama 2 / 3 & 7B--70B &
  7B: $\sim$8\,GB VRAM; 70B: $\sim$40\,GB &
  Baseline capable open weights; multi-GPU at 70B \\
\addlinespace
Mistral 7B & 7B &
  $\sim$7B-class &
  GQA + sliding-window attention; matches Llama 2 13B at lower cost \\
\addlinespace
Phi-3-mini & 3.8B & laptop GPU &
  Curated high-quality data; strong reasoning-per-param; edge / branch \\
\bottomrule
\end{tabular}

\medskip
\normalsize
\begin{alertblock}{Licences carry compliance obligations}
\textbf{Llama 2 Community License}: restricts orgs $>$700M MAU; derivative names must
include ``Meta Llama''.\quad\textbf{Mistral}: Apache 2.0 --- no such restriction.
Review licence terms before production.
\end{alertblock}
\end{frame}

\begin{frame}{Infrastructure: On-Premise vs.\ Private Cloud}
\begin{columns}[T]
\begin{column}{0.48\textwidth}
\textbf{On-premise}
\begin{itemize}
  \item[+] No data leaves the perimeter (prompts, completions, metadata)
  \item[--] GPU capex; ops burden; less elastic
\end{itemize}
\medskip
\textbf{Private cloud}
\begin{itemize}
  \item Single-tenant / VPC, dedicated HW
  \item AWS Outposts, Azure Gov, GCP Assured Workloads
  \item[!] DPA must satisfy GDPR Art.\ 28 \& DORA Art.\ 30
\end{itemize}
\end{column}
\begin{column}{0.52\textwidth}
\begin{exampleblock}{Routing for a commercial bank}
A lightweight classifier sorts each request:
\begin{enumerate}
  \item \textbf{Public}: no client IDs $\to$ cloud API
  \item \textbf{Internal}: internal analysis, no PII $\to$ private-cloud single-tenant
  \item \textbf{Client}: names / accounts / txns $\to$ on-prem, \emph{no external network}
\end{enumerate}
The classifier trains on internal data, calls no external API. Every routing
decision is logged (timestamp + tier).
\end{exampleblock}
\end{column}
\end{columns}
\end{frame}

\begin{frame}{Performance Trade-offs and Quantisation}
\textbf{The cost of going local is capability}: best open weights (2025) $\approx$
GPT-3.5-class; lag GPT-4 on complex reasoning and broad factual recall.

\medskip
\begin{underhood}
\textbf{Quantisation} lowers weight precision (16/32-bit $\to$ 4-bit).
\begin{itemize}
  \item \textbf{GPTQ} (Frantar et al., 2022): block-wise 2nd-order reconstruction
        using the Hessian; hours; needs a GPU.
  \item \textbf{AWQ} (Lin et al., 2023): $\sim$1\% of weights are \emph{salient}
        (large activations). Protect them with a per-channel scale; quantise the rest
        uniformly. Minutes; no GPU needed.
\end{itemize}
Similar final perplexity. \textbf{The quant step uses only a small \emph{public}
calibration set --- no client data, no added privacy risk.}
\end{underhood}

\medskip
\begin{block}{Memory at 4-bit}
7B $\approx$ \textbf{3.5\,GB} (less than a video game) $\cdot$ 13B $\approx$ 6.5\,GB
--- fits a high-end consumer GPU. A fine-tuned model can be quantised on-prem in $<$1 hr.
\end{block}
\end{frame}

% =====================================================================
\section{Text Anonymisation and De-identification}
% =====================================================================

\begin{frame}{When Local Deployment Is Not Enough}
\begin{bigpicture}
When inference must pass through external systems, the alternative is to remove or
obscure sensitive information \emph{before} the text reaches the model. The
technique --- de-identification --- comes from clinical informatics; finance swaps
diagnoses for client names, account numbers, LEIs, and trading positions.
\end{bigpicture}

\medskip
\textbf{The challenge}: sensitive data sits in \emph{prose}, not labelled fields ---
a loan officer's narrative, call notes, an analyst's draft. Removing a name without
context breaks intelligibility; replacing it must be \emph{consistent} so
co-references stay coherent.
\end{frame}

\begin{frame}{NER for Sensitive Financial Data}
Modern NER = pre-trained transformer encoders, fine-tuned; F1 $>$ 0.90 on standard
entity types (person, org, location, date, money) in well-edited English.

\medskip
\textbf{Finance adds entity types general NER misses} --- detected by regex or a
fine-tuned NER layer:
\begin{itemize}
  \item Account IDs: IBAN, SWIFT / BIC, account numbers
  \item Securities: ISIN (12-char), CUSIP, tickers
  \item Legal entities: LEI (20-char), registration numbers
  \item Regulatory: passport, national ID, tax ID
\end{itemize}

\medskip
\begin{block}{Production pipeline = two layers (spaCy, Honnibal \& Montani 2017)}
\textbf{Learned model} for prose entities (names, orgs, locations) $+$
\textbf{rule-based matcher} for structured identifiers (IBAN, ISIN, LEI).
\end{block}
\end{frame}

\begin{frame}{Worked NER Example}
\footnotesize
\emph{``The credit committee reviewed the application of Antoine Dupont (NIF:
12345678A) for a mortgage of EUR~320{,}000 on the property at 4 Avenue des
Champs-\'Elys\'ees, Paris 75008, with IBAN FR76 3000 6000 0112 3456 7890 189.''}

\medskip
\begin{tabular}{lll}
\toprule
Span & Type & Method \\
\midrule
Antoine Dupont & PERSON & Learned model \\
12345678A & TAX\_ID & Rule (NIF pattern) \\
EUR 320,000 & MONEY & Learned model \\
4 Avenue des Champs-\'Elys\'ees, Paris 75008 & ADDRESS & Learned model \\
FR76 3000 6000 0112 3456 7890 189 & IBAN & Rule (IBAN pattern) \\
\bottomrule
\end{tabular}

\medskip
\normalsize
\textbf{After masking:} \emph{``\ldots application of [PERSON\_001] (NIF:
[TAX\_ID\_001]) for a mortgage of [MONEY\_001] on the property at [ADDRESS\_001],
with IBAN [IBAN\_001].''} --- interpretable for summarisation, contains no PII.
\end{frame}

\begin{frame}{Three Anonymisation Strategies + $k$-Anonymity}
\begin{tabular}{p{2.2cm}p{4.0cm}p{4.0cm}}
\toprule
\textbf{Strategy} & \textbf{What it does} & \textbf{Trade-off} \\
\midrule
Masking & Typed placeholder \texttt{[PERSON\_001]} & Lossless structure; discards value \\
\addlinespace
Generalisation & Less specific value (75008 $\to$ Paris; EUR\,320k $\to$ EUR\,250--500k) & Keeps statistical content \\
\addlinespace
Suppression & Remove span entirely & Use when type itself is sensitive; can break coherence \\
\bottomrule
\end{tabular}

\medskip
\begin{block}{$k$-Anonymity (Def.\ 15.3, Sweeney 2002)}
$\mathcal{D}$ satisfies $k$-anonymity on quasi-identifiers $\mathcal{Q}$ if every
record is indistinguishable from $\geq k-1$ others when restricted to $\mathcal{Q}$.
For text: any output could plausibly come from $\geq k$ individuals. Masking
trivially gives large $k$ --- \emph{but} surrounding context may still re-identify.
\end{block}
\end{frame}

\begin{frame}{Pseudonymisation, Tokenisation, Re-identification}
\textbf{Pseudonymisation}: replace each entity with a realistic \emph{fake}
surrogate (``Marie Leclerc'' for ``Antoine Dupont''), consistent across the document.
\begin{itemize}
  \item[+] Reads natural --- LLM does not detect placeholders and degrade to generic output
  \item[--] The mapping table is a \textbf{key}: GDPR Recital 26 treats pseudonymised
        data as \emph{protected, not anonymised}. Hold the key separately,
        access-controlled, encrypted at rest.
\end{itemize}

\medskip
\textbf{Format-preserving tokenisation}: replace with a same-format token (valid-but-fake
IBAN). Standard in PCI-DSS; now in LLM pipelines --- data \emph{looks} real without being real.

\medskip
\begin{alertblock}{Contextual re-identification}
Even with all entities removed, a unique combination of dates, amount-ranges, and
product types may match one client. Risk $\propto$ adversary background knowledge
$\mathcal{K}$, which \emph{grows over time}. $\Rightarrow$ re-assess anonymisation
adequacy on a periodic cycle, especially for long-lived archives.
\end{alertblock}
\end{frame}

% =====================================================================
\section{Privacy-Preserving Training Techniques}
% =====================================================================

\begin{frame}{Bounding What the Model Memorises}
\begin{bigpicture}
Anonymisation guards the \emph{inference} channel. This section guards the
\emph{training} channel: limiting how much sensitive information becomes permanently
embedded during fine-tuning. Two complementary techniques ---
\textbf{differential privacy} bounds per-example contribution; \textbf{federated
learning} keeps data in place.
\end{bigpicture}
\end{frame}

\begin{frame}{Differential Privacy: The Definition}
\begin{block}{$(\varepsilon,\delta)$-Differential Privacy (Def.\ 15.4, Dwork \& Roth 2014)}
A randomised $\mathcal{M}$ is $(\varepsilon,\delta)$-DP if for all adjacent datasets
$D, D'$ (differing in one record) and all output sets $S$:
\[
\Pr[\mathcal{M}(D) \in S] \leq e^{\varepsilon}\,\Pr[\mathcal{M}(D') \in S] + \delta
\]
\end{block}
\begin{itemize}
  \item $\varepsilon \geq 0$ = \emph{privacy budget} (smaller = more private)
  \item $\delta \geq 0$ = \emph{failure probability}; pure DP has $\delta = 0$
\end{itemize}

\medskip
\textbf{Intuition.} An adversary observing the output cannot confidently tell whether
any record was in the input. Log-odds of distinguishing $D$ from $D'$ are bounded by
$\varepsilon$ (exactly, with probability $1-\delta$).
\end{frame}

\begin{frame}{The Gaussian Mechanism and DP-SGD}
\begin{block}{Gaussian Mechanism (Thm.\ 15.1)}
For $f$ with $\ell_2$-sensitivity $\Delta_2 f = \max_{D,D'}\|f(D)-f(D')\|_2$:
\[
\mathcal{M}(D) = f(D) + \mathcal{N}(0,\sigma^2 I_d),\quad
\sigma = \frac{\Delta_2 f \cdot \sqrt{2\ln(1.25/\delta)}}{\varepsilon}
\]
is $(\varepsilon,\delta)$-DP for $\varepsilon,\delta \in (0,1)$.
\end{block}

\medskip
\begin{underhood}
\textbf{DP-SGD} (Abadi et al., 2016) --- three changes to SGD:
\begin{enumerate}
  \item \emph{Per-example gradients} $g_i$ (not minibatch sum).
  \item \emph{Clip}: $\hat{g}_i = g_i / \max(1, \|g_i\|_2/C)$ --- bounds sensitivity to $C$.
  \item \emph{Noise}: add $\mathcal{N}(0, \sigma^2 C^2 I)$ to the clipped sum before the update.
\end{enumerate}
Budget tracked by the \textbf{moments accountant} / R\'enyi-DP --- tighter than basic
composition.
\end{underhood}
\end{frame}

\begin{frame}{Privacy--Utility Trade-off and Composition}
\begin{columns}[T]
\begin{column}{0.55\textwidth}
\textbf{The cost of DP-SGD is utility.} Noise slows convergence.
\begin{itemize}
  \item $\varepsilon = 1$: strong privacy, larger utility cost
  \item $\varepsilon = 10$: weak privacy, small cost
\end{itemize}
Accuracy rises with $\varepsilon$ and \emph{saturates} toward the non-private baseline.

\medskip
\textbf{LoRA / adapters pair well with DP-SGD}: fewer trainable params $\Rightarrow$
lower sensitivity $\Rightarrow$ less noise.
\end{column}
\begin{column}{0.45\textwidth}
\begin{block}{Composition (Prop.\ 15.1)}
Sequential: $(\varepsilon_1{+}\varepsilon_2,\,\delta_1{+}\delta_2)$.\\[4pt]
Advanced, $k$ steps:
\[
\varepsilon' \approx \varepsilon\sqrt{2k\ln(1/\delta')}
\]
\end{block}
\end{column}
\end{columns}

\medskip
\begin{alertblock}{Operational implication}
Every DP fine-tune \emph{spends} budget. Plan cumulative loss across runs; once
exhausted, further training violates the guarantee --- requiring a fresh model or a
new budget allocation.
\end{alertblock}
\end{frame}

\begin{frame}{Federated Learning: FedAvg}
\textbf{Different threat: data centralisation.} Data-residency rules may forbid
moving client data across a jurisdiction --- even within one institution.

\medskip
\begin{block}{Federated Averaging (Def.\ 15.5, McMahan et al.\ 2017)}
$K$ clients hold $D_k$ ($n_k$ examples, $n=\sum n_k$). Minimise
$F(\theta)=\sum_k \frac{n_k}{n}F_k(\theta)$. Each round $t$: selected clients init
with $\theta^t$, run $E$ local steps, return $\Delta_k^t$. Server aggregates:
\[
\theta^{t+1} = \theta^t + \sum_{k \in S_t}\frac{n_k}{\sum_{j\in S_t} n_j}\,\Delta_k^t
\]
\end{block}

\medskip
\textbf{Guarantee \& limit}: the aggregator sees only \emph{updates}, never raw data
$\Rightarrow$ data-residency satisfied. But updates leak via \textbf{gradient
inversion} --- combine with \emph{local} DP-SGD for both guarantees.

\medskip
\textbf{Finance use cases}: (1) multi-bank joint fraud / AML model without sharing
data; (2) cross-border improvement of a single global model.
\end{frame}

% =====================================================================
\section{Architectural Patterns for Privacy-Compliant Systems}
% =====================================================================

\begin{frame}{Privacy-Aware RAG}
\textbf{RAG leaks at three points}: ingestion (embed \& store), retrieval (query
embedding reaches the vector DB operator), inference (retrieved docs sit beside user data).

\medskip
\begin{block}{The systematic fix}
\begin{itemize}
  \item \textbf{Ingestion}: anonymise \emph{before} embedding; store only the
        anonymised version; keep the mapping in an access-controlled store.
  \item \textbf{Retrieval}: compute embeddings with a \emph{local} model --- query
        never leaves the network.
  \item \textbf{Inference}: scan retrieved chunks for residual PII before prompting.
  \item \textbf{Access control}: row-level security --- each chunk carries an ACL;
        retrieval is filtered by the user's clearance.
\end{itemize}
\end{block}

\medskip
\textbf{Why access control matters}: a junior analyst must not retrieve merger-target
docs via an indirect question to the assistant.
\end{frame}

\begin{frame}{Audit Trails and Encrypted Pipelines}
Regulators expect full reconstruction of any client-affecting decision: \emph{who}
asked, \emph{what} was retrieved, \emph{what} was asked / answered, \emph{what} action followed.

\medskip
\begin{exampleblock}{Privacy-aware audit logging}
Logging full prompt + completion is the most complete record --- but creates a
\emph{secondary} PII store under GDPR retention + DORA resilience. Better: log a
\textbf{hash} of the prompt (integrity, not content) plus structured metadata ---
user ID, clearance, document IDs, model version, response time, outcome --- to an
\emph{immutable} store.
\end{exampleblock}

\medskip
\textbf{Encryption end-to-end}: at rest, in transit, ideally during inference.
\begin{itemize}
  \item \emph{Homomorphic encryption}: compute on ciphertext --- infeasible overhead
        for transformer inference today.
  \item \emph{Confidential computing} (TEE: Intel TDX, AMD SEV): practical;
        blocks even the cloud operator during processing.
\end{itemize}
\end{frame}

% =====================================================================
\section{Evaluation and Compliance}
% =====================================================================

\begin{frame}{Measuring Privacy Leakage Empirically}
\begin{bigpicture}
\textbf{Verify privacy --- do not take it on faith from the training algorithm.}
Three standard quantitative tests, run after every fine-tune.
\end{bigpicture}

\medskip
\begin{enumerate}
  \item \textbf{Membership inference AUC}: hold out fine-tuning docs as negatives; run
        the likelihood-ratio attack; AUC of distinguishing members from non-members.
        $0.5$ = perfect privacy; $1.0$ = total failure. \textbf{$<0.6$ commonly
        acceptable.}
  \item \textbf{Verbatim extraction rate}: insert ``canary'' sensitive strings; measure
        the fraction recovered. DP fine-tuning $\Rightarrow$ near-zero.
  \item \textbf{Property inference accuracy}: train a meta-classifier on shadow models
        to infer an aggregate property; compare to a random baseline.
\end{enumerate}
\end{frame}

\begin{frame}{Compliance Checklist (Selected Rows)}
\footnotesize
\begin{tabular}{p{2.7cm}p{1.9cm}p{2.7cm}p{2.6cm}}
\toprule
\textbf{Obligation} & \textbf{Source} & \textbf{Mitigation} & \textbf{Evidence} \\
\midrule
Restrict automated decisions & GDPR Art.\ 22 & Human-in-the-loop review & Workflow docs \\
Right to erasure & GDPR Art.\ 17 & Exclude subject from fine-tune; use retrieval & Architecture docs \\
Data minimisation & GDPR Art.\ 5 & Anonymise before embed; log hashes & Pipeline design \\
DPIA (high risk) & GDPR Art.\ 35 & Conduct before deployment & Signed DPIA \\
ICT third-party contract & DORA Art.\ 30 & DORA-compliant DPA & Signed contract \\
Record-keeping & MiFID II & Immutable audit trail & Log samples \\
Model risk mgmt & SR 11-7 / ECB & Model card, benchmarks, oversight & Model docs \\
\bottomrule
\end{tabular}

\smallskip
A \textbf{DPIA} (GDPR Art.\ 35) must precede any LLM deployment processing personal
financial data --- and must address memorisation risk, the third-party data-sharing
model, and the erasure mechanism.
\end{frame}

\begin{frame}{Lecture 15 Summary: Mitigations by Threat}
\begin{block}{Five strategies, complementary --- apply \emph{all} in production}
\small
\begin{tabular}{p{5.0cm}p{5.0cm}}
\toprule
\textbf{Threat} & \textbf{Mitigation} \\
\midrule
Inference-channel exfiltration & Local / private-cloud deployment \\
Input disclosure via prompt & Text anonymisation / de-identification \\
Training-channel memorisation & DP-SGD during fine-tuning \\
Data-residency violation & Federated learning \\
Prompt injection / access failure & Privacy-aware RAG + audit logging \\
\bottomrule
\end{tabular}
\end{block}

\medskip
\normalsize
\begin{alertblock}{The residual gap}
\textbf{Right to erasure (GDPR Art.\ 17)}: excluding a subject from future
fine-tuning is easy; \emph{removing} already-memorised information from a deployed
model --- \textbf{machine unlearning} --- is an open research problem, not a solved
control. Re-measure residual exposure (MI-AUC, canary rate) on a fixed cadence.
\end{alertblock}
\end{frame}

% =====================================================================
\appendixstart{Extended Derivations, Internals, and Edge Cases}
% =====================================================================

\begin{frame}[noframenumbering]{Appendix A --- Advanced Composition in Full}
\begin{block}{Proposition 15.1 (full statement)}
For $k$ mechanisms each $(\varepsilon,\delta)$-DP, the composite is
$(\varepsilon', k\delta+\delta')$-DP for any $\delta'>0$ with
\[
\varepsilon' = \varepsilon\sqrt{2k\ln(1/\delta')} + k\varepsilon(e^{\varepsilon}-1).
\]
For small $\varepsilon$ ($e^{\varepsilon}-1\approx\varepsilon$) the second term is
$O(k\varepsilon^2)$, often negligible vs.\ the first ($\Rightarrow$ the cited
$\varepsilon'\approx\varepsilon\sqrt{2k\ln(1/\delta')}$).
\end{block}

\medskip
\textbf{Proof sketch.} Basic composition uses post-processing immunity: any function
of a DP output keeps the guarantee. Chain the two per-mechanism bounds via a union
argument $\Rightarrow (\varepsilon_1{+}\varepsilon_2,\,\delta_1{+}\delta_2)$. The
advanced bound follows from the moments accountant tracking the R\'enyi divergence
$D_\alpha(P\|Q)=\frac{1}{\alpha-1}\log\sum_x P(x)^\alpha Q(x)^{1-\alpha}$ between
output distributions on $D, D'$ rather than the multiplicative distance.
\end{frame}

\begin{frame}[noframenumbering]{Appendix B --- AWQ vs.\ GPTQ Internals}
\begin{underhood}
The 16$\to$4-bit error is \emph{not uniform}. \textbf{AWQ}'s insight: $\sim$1\% of
weights are salient (large activation magnitudes); aggressive quantisation of these
causes large output errors. AWQ finds them from activation statistics on a small
calibration set, then applies a per-channel \emph{scale} transformation that protects
them without raising bit width; the other 99\% go uniform 4-bit with negligible loss.

\medskip
\textbf{GPTQ} instead solves a block-wise reconstruction per weight matrix,
minimising quantisation error given the Hessian of outputs w.r.t.\ weights.

\medskip
Similar final perplexity; AWQ is faster (minutes vs.\ hours, 7B) and needs no GPU.
\end{underhood}
\end{frame}

\begin{frame}[noframenumbering]{Appendix C --- Re-identification as Inference}
Model re-identification as an inference problem: the adversary observes the
anonymised document $\hat{d}$ and uses prior knowledge $\mathcal{K}$ to form
\[
\mathcal{P}(\text{individual} \mid \hat{d}, \mathcal{K}).
\]
\begin{itemize}
  \item Posterior concentrated on one individual $\Rightarrow$ re-identification trivial.
  \item Posterior uniform over a large set $\Rightarrow$ re-identification infeasible.
\end{itemize}

\medskip
\textbf{Assessment by simulation}: apply the pipeline, then attempt re-identification
using the institution's \emph{own} databases as $\mathcal{K}$. Success $\Rightarrow$
pipeline insufficient; strengthen generalisation / suppression. Failure on all test
cases $\Rightarrow$ adequate \emph{for that threat model}. Point-in-time only ---
$\mathcal{K}$ grows; re-review on a cycle.
\end{frame}

\begin{frame}[noframenumbering]{Appendix D --- Privacy-Aware RAG Walkthrough}
\textbf{Credit analyst query:} ``Typical collateral structure for technology-sector LBO loans?''
\begin{enumerate}
  \item Query embedded \emph{locally} by a self-hosted model.
  \item Top-5 chunks retrieved, filtered to the analyst's authorised set (policy docs,
        sector reports, anonymised historical loan files).
  \item Chunks scanned for PII --- none (pre-anonymised at ingestion).
  \item Local LLM generates, citing retrieved content.
  \item Logged: timestamp, user ID, \emph{query hash} (not the query), document IDs,
        model version $\to$ immutable audit store.
\end{enumerate}

\medskip
\textbf{Result}: no client-identifying information enters the prompt; no data crosses
an external boundary.
\end{frame}

\end{document}
